\documentclass[12pt, letterpaper]{article}
\usepackage[margin=1in]{geometry}
\usepackage{amsmath, amssymb}
\usepackage{booktabs}
\usepackage{array}
\usepackage{xcolor}
\usepackage{hyperref}
\usepackage{enumitem}
\usepackage{titlesec}
\usepackage{fancyhdr}
\usepackage{parskip}
\usepackage{setspace}
\usepackage{float}
\usepackage{listings}
\usepackage[T1]{fontenc}
\usepackage[utf8]{inputenc}

\definecolor{burnt}{RGB}{191,87,0}
\definecolor{lightgray}{RGB}{245,245,245}
\definecolor{codeblue}{RGB}{0,90,160}

\hypersetup{colorlinks=true, linkcolor=burnt, urlcolor=burnt, citecolor=burnt}

\titleformat{\section}{\large\bfseries\color{burnt}}{}{0em}{}[\color{burnt}\titlerule]
\titleformat{\subsection}{\normalsize\bfseries\color{burnt!80}}{}{0em}{}
\titleformat{\subsubsection}{\normalsize\itshape}{}{0em}{}

\lstset{
  backgroundcolor=\color{lightgray},
  basicstyle=\ttfamily\small,
  keywordstyle=\color{codeblue}\bfseries,
  commentstyle=\color{gray},
  stringstyle=\color{burnt},
  breaklines=true,
  frame=single,
  numbers=left,
  numberstyle=\tiny\color{gray},
  language=R
}

\pagestyle{fancy}
\fancyhf{}
\rhead{\textcolor{burnt}{\small Big Data \& Data Mining}}
\lhead{\textcolor{gray}{\small UT Austin}}
\rfoot{\textcolor{gray}{\small Page \thepage}}
\lfoot{\textcolor{gray}{\small sb74887 $\cdot$ ha26947}}
\renewcommand{\headrulewidth}{0.4pt}
\renewcommand{\footrulewidth}{0.4pt}

\begin{document}

\begin{titlepage}
  \centering
  \vspace*{2cm}
  {\color{burnt}\rule{\linewidth}{2pt}}\\[0.5cm]
  {\LARGE\bfseries Credit Card Customer Behavior Analysis}\\[0.3cm]
  {\large Hypothesis Testing \& K-Means Clustering}\\[0.3cm]
  {\color{burnt}\rule{\linewidth}{2pt}}
  \vspace{1.5cm}
  {\large
    \begin{tabular}{rl}
      \textbf{Course:}      & Big Data \& Data Mining \\[0.3em]
      \textbf{Instructor:}  & Dr.\ Kia Teymourian \\[0.3em]
      \textbf{Institution:} & The University of Texas at Austin \\[0.3em]
      \textbf{Authors:}     & Hannia Ashley Alvarado Galv\'{a}n \quad \texttt{ha26947} \\
                            & Santiago Basald\'{u}a Ram\'{i}rez \quad \texttt{sb74887} \\[0.3em]
      \textbf{Date:}        & \today \\
    \end{tabular}
  }
  \vspace{1.5cm}
  \begin{abstract}
    \noindent
    This report documents the methodology, analysis, and results of a data analytics
    project using a real-world credit card dataset containing 8{,}950 customer records
    across 18 behavioral and financial variables. The project has two parts.
    \textbf{Part 1} applies formal hypothesis testing to evaluate two claims:
    (1) that customers who frequently use cash advances carry significantly higher
    account balances (Welch's two-sample $t$-test), and (2) that account balance
    is a strong predictor of credit limit (Pearson correlation --- debunked).
    \textbf{Part 2} applies K-Means clustering to discover five distinct and
    interpretable customer behavioral profiles. All analysis was performed in R
    using \texttt{dplyr} for data wrangling and base R statistical functions.
  \end{abstract}
  \vfill
  {\small\textit{Dataset: CC GENERAL --- Kaggle Credit Card Customer Segmentation}}
\end{titlepage}

\tableofcontents
\newpage
\onehalfspacing

% ============================================================================
\section{Introduction}
% ============================================================================

Understanding how customers use credit cards is a fundamental problem in financial
data science. Card issuers need to identify risky behaviors early, set appropriate
credit limits, and segment customers to offer personalized services.

This project addresses the central research question:

\begin{center}
  \textit{\large ``Do distinct patterns of credit card usage carry statistically
  measurable differences, and can we identify natural customer profiles through
  clustering?''}
\end{center}

The analysis is structured in two complementary parts:

\begin{enumerate}
  \item \textbf{Hypothesis Testing} --- Confirm or debunk specific claims about
    customer financial behavior using formal statistical tests.
  \item \textbf{K-Means Clustering} --- Apply unsupervised machine learning to
    discover naturally occurring customer segments.
\end{enumerate}

% ============================================================================
\section{Dataset Description}
% ============================================================================

The dataset is from Kaggle's \textit{Credit Card Customer Segmentation} collection.
It contains \textbf{8{,}950 customer records} with \textbf{18 variables} covering
spending behavior, balance management, and credit utilization over six months.

\subsection{Variable Dictionary}

\begin{table}[H]
  \centering
  \caption{Official Variable Definitions}
  \renewcommand{\arraystretch}{1.3}
  \begin{tabular}{>{\ttfamily\small}p{5.5cm} p{8cm}}
    \toprule
    \normalfont\textbf{Variable} & \textbf{Official Definition} \\
    \midrule
    CUST\_ID                       & Identification of Credit Card holder (Categorical) \\
    BALANCE                        & Balance amount left in account to make purchases \\
    BALANCE\_FREQUENCY             & How frequently Balance is updated; score 0--1 \\
    PURCHASES                      & Amount of purchases made from account \\
    ONEOFF\_PURCHASES              & Maximum purchase amount done in one-go \\
    INSTALLMENTS\_PURCHASES        & Amount of purchases done in installment \\
    CASH\_ADVANCE                  & Cash in advance given by the user \\
    PURCHASES\_FREQUENCY           & How frequently purchases are made; score 0--1 \\
    ONEOFFPURCHASES\_FREQUENCY     & How frequently one-go purchases happen; score 0--1 \\
    PURCHASES\_INSTALLMENTS\_FREQ  & How frequently installment purchases are done; 0--1 \\
    CASH\_ADVANCE\_FREQUENCY       & How frequently cash in advance is paid; score 0--1 \\
    CASH\_ADVANCE\_TRX             & Number of transactions made with Cash in Advanced \\
    PURCHASES\_TRX                 & Number of purchase transactions made \\
    CREDIT\_LIMIT                  & Limit of Credit Card for user \\
    PAYMENTS                       & Amount of Payment done by user \\
    MINIMUM\_PAYMENTS              & Minimum amount of payments made by user \\
    PRCFULLPAYMENT                 & Percent of full payment paid by user \\
    TENURE                         & Tenure of credit card service for user \\
    \bottomrule
  \end{tabular}
\end{table}

\subsection{Note on BALANCE Interpretation}

Despite the official description saying ``balance left to make purchases,''
\texttt{BALANCE} in this dataset represents the \textbf{outstanding debt balance}
(what the customer owes). This is confirmed empirically: the correlation between
\texttt{CASH\_ADVANCE} and \texttt{BALANCE} is positive ($r \approx 0.50$).
If BALANCE were available credit, drawing cash advances would \textit{decrease} it.
Since it \textit{increases}, BALANCE must represent accumulated debt.

% ============================================================================
\section{Data Wrangling with dplyr}
% ============================================================================

All data wrangling was performed using the \texttt{dplyr} package in R.

\subsection{Cleaning}

Rows with missing values in \texttt{CREDIT\_LIMIT} and \texttt{MINIMUM\_PAYMENTS}
were removed. After cleaning, \textbf{8{,}636 records} remained (314 removed, $<$4\%).

\subsection{New Variables}

\begin{itemize}
  \item \textbf{\texttt{CA\_Group}}: Binary label: \textit{Heavy CA} if
    \texttt{CASH\_ADVANCE\_FREQUENCY} $\geq 0.25$, else \textit{Light CA}.
    Threshold = 75th percentile of the frequency distribution.
  \item \textbf{\texttt{LOG\_BALANCE}}: $\log(1 + \texttt{BALANCE})$ to handle
    right-skewness (skewness $\approx 2.39$) for $t$-test normality requirements.
\end{itemize}

\begin{lstlisting}[caption={Data Wrangling Pipeline}]
cc_clean <- cc %>%
  filter(!is.na(CREDIT_LIMIT), !is.na(MINIMUM_PAYMENTS)) %>%
  mutate(
    CA_Group    = if_else(CASH_ADVANCE_FREQUENCY >= 0.25,
                          "Heavy CA", "Light CA"),
    LOG_BALANCE = log1p(BALANCE)
  )
\end{lstlisting}

% ============================================================================
\section{Part 1: Hypothesis Testing}
% ============================================================================

All tests use significance level $\alpha = 0.05$.

% ---------------------------------------------------------------------------
\subsection{H1: Cash-Advance Users Carry Higher Balances}
% ---------------------------------------------------------------------------

\subsubsection{Formal Statement}

\[
  H_0: \mu_{\text{Heavy CA}} = \mu_{\text{Light CA}}
  \qquad
  H_1: \mu_{\text{Heavy CA}} > \mu_{\text{Light CA}}
  \quad\text{(one-tailed)}
\]

Where $\mu$ denotes the mean account balance of each group.

\subsubsection{Why Welch's $t$-Test?}

The standard two-sample $t$-test assumes equal variances. Heavy CA users have a
wider spread of balances. Welch's modification uses separate variance estimates
per group and adjusted degrees of freedom --- the safe default for unequal variances.

\begin{lstlisting}[caption={H1 Test}]
t_result_h1 <- t.test(heavy, light,
                       alternative = "greater",
                       var.equal   = FALSE)
\end{lstlisting}

\subsubsection{Results}

\begin{table}[H]
  \centering
  \caption{H1 Results}
  \begin{tabular}{lc}
    \toprule
    \textbf{Statistic} & \textbf{Value} \\
    \midrule
    $t$-statistic            & $\approx 28$ \\
    Degrees of freedom       & $\approx 2{,}889$ \\
    $p$-value                & $< 2.2\times10^{-16}$ \\
    95\% CI lower bound      & $>\$1{,}000$ \\
    Decision                 & \textbf{REJECT $H_0$} \\
    \bottomrule
  \end{tabular}
\end{table}

\textbf{Conclusion:} Heavy CA users carry significantly higher balances.
\texttt{CASH\_ADVANCE\_FREQUENCY} is a reliable credit-risk signal.

% ---------------------------------------------------------------------------
\subsection{H2: Balance Strongly Predicts Credit Limit (Debunk)}
% ---------------------------------------------------------------------------

\subsubsection{The Intuitive Claim}

``Customers with higher balances must have higher credit limits.''

\subsubsection{Formal Statement}

\[
  H_0: \rho = 0
  \qquad
  H_1: \rho \neq 0 \text{ (strong positive correlation)}
\]

\subsubsection{Debunk Strategy}

\begin{enumerate}
  \item \textbf{Moderate $r$} --- $r^2$ must be large to call it ``strong.''
  \item \textbf{Outliers} --- IQR rule: $>$666 balance outliers distort the pattern.
  \item \textbf{Right-skewed distribution} --- Skewness $\approx 2.39$ violates
    the normality assumption of Pearson correlation.
\end{enumerate}

\begin{lstlisting}[caption={H2 Test}]
cor_h2 <- cor.test(cc_clean$BALANCE, cc_clean$CREDIT_LIMIT,
                   method = "pearson")
\end{lstlisting}

\subsubsection{Results}

\begin{table}[H]
  \centering
  \caption{H2 Results}
  \begin{tabular}{lc}
    \toprule
    \textbf{Statistic} & \textbf{Value} \\
    \midrule
    $r$ (Pearson)              & $\approx 0.54$ \\
    $r^2$ (Variance explained) & $\approx 29\%$ \\
    $p$-value                  & $< 2.2\times10^{-16}$ \\
    BALANCE outliers (IQR)     & $>666$ customers \\
    Decision                   & \textbf{DEBUNKED} \\
    \bottomrule
  \end{tabular}
\end{table}

\subsubsection{Key Lesson: Significance $\neq$ Practical Strength}

With $n = 8{,}636$, even moderate correlations produce near-zero $p$-values.
BALANCE explains only \textbf{29\%} of CREDIT\_LIMIT variance. The other 71\%
depends on income, employment history, and credit score --- factors not in this
dataset. Banks set limits \textit{prospectively} based on repayment capacity,
not reactively based on current debt.

\textbf{Conclusion: Debunked.} BALANCE is not a strong predictor of CREDIT\_LIMIT.

% ============================================================================
\section{Part 2: K-Means Clustering}
% ============================================================================

\subsection{Motivation}

Hypothesis testing tested specific claims. Clustering asks an open question:
\textit{``Can we find natural customer groups in the data?''}
K-Means is \textbf{unsupervised}: no predefined labels are required.

\subsection{Variable Selection}

Eight variables were chosen for low inter-correlation and behavioral relevance:
BALANCE, PURCHASES, CASH\_ADVANCE, CREDIT\_LIMIT, PAYMENTS,
PURCHASES\_FREQUENCY, CASH\_ADVANCE\_FREQUENCY, PRCFULLPAYMENT.

High-redundancy pairs (e.g., PURCHASES and ONEOFF\_PURCHASES, $r\approx0.92$)
were excluded.

\subsection{Standardization}

K-Means uses Euclidean distance. Dollar-valued variables (BALANCE: 0--19{,}000)
would dominate frequency variables (0--1) without standardization.
All variables were rescaled to mean 0, standard deviation 1 using \texttt{scale()}.

\subsection{Choosing K = 5}

$K \in \{2,\ldots,8\}$ were evaluated via the elbow method (within-cluster inertia).
$K=5$ was selected: it falls at the elbow and produces five \textbf{interpretable}
business profiles.

\begin{lstlisting}[caption={K-Means Setup}]
X_scaled <- scale(cc_clust)
set.seed(123)
km <- kmeans(X_scaled, centers = 5, nstart = 25, iter.max = 300)
\end{lstlisting}

\texttt{nstart = 25} runs 25 random initializations; the best result is kept.

\subsection{Cluster Profiles}

\begin{table}[H]
  \centering
  \caption{Five Customer Profiles}
  \renewcommand{\arraystretch}{1.35}
  \begin{tabular}{p{4.5cm} c p{6.5cm}}
    \toprule
    \textbf{Cluster} & \textbf{\%} & \textbf{Key Characteristics} \\
    \midrule
    Low-Activity Customers  & $\sim$34\% & Very low purchases and frequency; minimal card use \\
    Regular Purchasers      & $\sim$26\% & Frequent purchases; very low cash-advance use \\
    Full-Payment Customers  & $\sim$18\% & Consistently pay full balance; lowest debt \\
    High CA Users           & $\sim$20\% & High balance driven by cash advances; highest risk \\
    High-Value Customers    & $\sim$1.3\% & Extremely high spending and credit limit ($n\approx112$) \\
    \bottomrule
  \end{tabular}
\end{table}

\begin{description}
  \item[Low-Activity] Largest segment. Likely keep the card for emergencies or
    to maintain an open credit line for their credit score.
  \item[Regular Purchasers] Active, responsible users. Purchase frequently,
    avoid cash advances, pay regularly.
  \item[Full-Payment] Most financially disciplined. Pay entire balance most months.
    Corroborates $r(\texttt{PRCFULLPAYMENT}, \texttt{BALANCE})\approx-0.32$.
  \item[High CA Users] Highest-risk group. Compounding interest from cash advances
    drives up their balance (independently corroborates H1).
  \item[High-Value] Tiny elite ($n\approx112$, 1.3\%). Likely high-income
    individuals or business owners with outsized financial impact.
\end{description}

% ============================================================================
\section{Key Correlations}
% ============================================================================

\begin{table}[H]
  \centering
  \caption{Most Important Bivariate Correlations}
  \begin{tabular}{lcc}
    \toprule
    \textbf{Variable Pair} & \textbf{$r$} & \textbf{Interpretation} \\
    \midrule
    PURCHASES $\leftrightarrow$ PAYMENTS          & $+0.60$ & Spend more, pay more \\
    CASH\_ADVANCE $\leftrightarrow$ BALANCE       & $+0.50$ & CA use accumulates debt \\
    PRCFULLPAYMENT $\leftrightarrow$ BALANCE      & $-0.32$ & Full payers carry less debt \\
    PURCHASES $\leftrightarrow$ ONEOFF\_PURCHASES & $+0.92$ & Redundant (excluded from clustering) \\
    BALANCE $\leftrightarrow$ CREDIT\_LIMIT       & $+0.54$ & Moderate only (H2 tested) \\
    \bottomrule
  \end{tabular}
\end{table}

% ============================================================================
\section{Exploratory Data Analysis}
% ============================================================================

\begin{itemize}
  \item \textbf{8{,}950 customers}, 18 variables. All major financial variables
    are strongly right-skewed:
    \begin{itemize}
      \item PURCHASES: skewness $= 8.14$
      \item CASH\_ADVANCE: skewness $= 5.17$
      \item BALANCE: skewness $= 2.39$
    \end{itemize}
  \item Mean $\gg$ Median for most variables: a small group of high-activity
    customers pulls the average upward.
  \item Extreme values were \textbf{kept}: they represent real behavior and
    define the High-Value and High CA clusters.
\end{itemize}

% ============================================================================
\section{Conclusions}
% ============================================================================

\subsection{Part 1: Hypothesis Testing}

\begin{itemize}
  \item \textbf{H1 Confirmed} --- Heavy CA users carry significantly higher balances
    ($p \ll 0.05$, $t\approx28$). \texttt{CASH\_ADVANCE\_FREQUENCY} is a reliable
    financial risk signal.
  \item \textbf{H2 Debunked} --- Despite a near-zero $p$-value, BALANCE explains
    only $\approx29\%$ of CREDIT\_LIMIT variance ($r\approx0.54$, $r^2\approx0.29$).
    Statistical significance does not imply practical predictive strength.
\end{itemize}

\subsection{Part 2: Clustering}

K-Means found \textbf{five interpretable customer profiles}. These are exploratory,
not official categories. The High CA Users cluster independently confirms H1 by
showing the highest average balance alongside maximum cash-advance activity.

\subsection{Key Methodological Lesson}

The most important takeaway: \textbf{significance $\neq$ practical relevance}.
With large $n$, trivially small effects become significant. Always report
$r^2$ or Cohen's $d$ alongside the $p$-value.

% ============================================================================
\section{References}
% ============================================================================

\begin{itemize}
  \item \textbf{Dataset:} CC GENERAL --- Kaggle.
    \url{https://www.kaggle.com/datasets/arjunbhasin2013/ccdata}
  \item \textbf{R Core Team} (2024). R: A language and environment for statistical
    computing. R Foundation for Statistical Computing.
  \item \textbf{Wickham, H. et al.} (2023). \textit{dplyr}: A Grammar of Data
    Manipulation. R package version 1.1.4.
  \item \textbf{Welch, B.~L.} (1947). The generalization of Student's problem.
    \textit{Biometrika}, 34(1--2), 28--35.
  \item \textbf{MacQueen, J.} (1967). Some methods for classification and analysis
    of multivariate observations. \textit{Proc. Fifth Berkeley Symposium}, 1, 281--297.
\end{itemize}

\vspace{2em}
\begin{center}
  \color{gray}\rule{0.5\linewidth}{0.4pt}\\[0.3em]
  {\small Big Data \& Data Mining --- Dr.\ Kia Teymourian ---
  The University of Texas at Austin}\\
  {\small \texttt{sb74887} $\cdot$ \texttt{ha26947}}
\end{center}

\end{document}
